\documentclass[11pt,a4paper]{article}
\usepackage[T1]{fontenc}
\usepackage{lmodern,amsmath,amssymb,amsthm,mathtools,booktabs,microtype}
\usepackage[margin=27mm]{geometry}
\usepackage[colorlinks=true,linkcolor=blue,citecolor=blue,urlcolor=blue]{hyperref}
\setlength{\emergencystretch}{2em}
\newtheorem{theorem}{Theorem}[section]
\newtheorem{lemma}[theorem]{Lemma}
\newtheorem{proposition}[theorem]{Proposition}
\newtheorem{corollary}[theorem]{Corollary}
\newcommand{\Q}{\mathbb Q}\newcommand{\Z}{\mathbb Z}
\DeclareMathOperator{\Tr}{Tr}\DeclareMathOperator{\Nm}{N}
\title{Norm preparation for general and conditional elliptic detectors}
\author{}\date{16 November 2026}
\begin{document}\maketitle
\begin{abstract}
We give an explicit degree-eight input polynomial and a compatible scalar
coefficient for the general elliptic detector. The input polynomial
reflects rational integrality, and the scalar is a norm from
$\Q(\sqrt{10657})$.
The scalar retains a uniformly positive odd valuation at the consumed
prime, square reduced denominator on integer inputs, and the required
local preparation at $2,3,17$. The same arithmetic premises and retained
coordinates give a twenty-five-variable family. Its prepared contact
weights can be chosen to be norms, so both associated quadratic norm
conics have rational points. The unchanged conditional thirteen-variable
sources admit the same simultaneous norm preparation, without a scalar
replacement. Realization with the prescribed detector traces remains a
separate condition; no smaller variable count is asserted.
\end{abstract}

\section{Statement and retained arithmetic scope}
We use the current general source of~\cite{UnitNorm}, with the unsquared
scale $a_1=1+\rho T^2$. Its retained arithmetic premises are those of
\cite{General,Dependencies}: the specified elliptic subgroup and index
height, the normalized modular realization giving the five-contact
height estimate, the prepared-input pole theorem, and the simultaneous
prime-pattern theorem with fixed avoidance and prescribed Gaussian
residues. The aggregate trace-height theorem is used in precisely the
scope established there. No additional rational $abc$ hypothesis is used
for this general family.

Put
\[
 \ell=76249,\qquad d_*=10657,\qquad F_* =\Q(\sqrt{d_*}),
\]
and let $S$ be the original exceptional set: the primes at most $76176$
together with $1875611,3828397,3982747$. In particular $d_*\in S$
and $\ell\notin S$. Define
\begin{equation}\label{eq:constants}
\begin{gathered}
 n=116576341082536192,\quad D_0=5951086353,\\
 K_0=25481059893697,\quad H_0=279986328=2^3 3^3 17\ell,\\
 A=97391047218196174706960881179098628096,\\
 B=152061501652330028986565424,\qquad C=59355302528062.
\end{gathered}
\end{equation}
Here the fixed coefficient $B$ is unrelated to the shifts in the
conditional thirteen-variable sources. Set
\begin{equation}\label{eq:input}
 U(X)=1+\frac{X(X-1)(X-2)(X-3)}{24},\qquad
 \beta_*(X)=AU(X)^2+BU(X)+C.
\end{equation}
With
\[
 r_4(t)=t(t-1)(t-2401)(32t+343),\qquad r_*(X)=r_4(\beta_*(X)),
\]
put
\begin{equation}\label{eq:scalar}
\begin{gathered}
 M_0=r_4(C)
 =397180638679647074411201284065117677360619724674885081954,\\
 \eta=34,\quad A_0=102\eta M_0,\quad
 C_L=A_0(\ell^{(L+1)/2}+1),\\
 \rho_{*,L}(X)=
 \frac{\eta^2\ell^L M_0r_*(X)}{(1+C_Lr_*(X))^2}
 \qquad(L\ge3\text{ odd}).
\end{gathered}
\end{equation}

\begin{theorem}\label{thm:main}
Under the retained arithmetic premises, substitution of
$\beta_*,\rho_{*,L}$ into the general source gives a computable family
of integer polynomials in twenty-five total variables. Every sufficiently
large fixed odd $L$ defines $\Z$ over $\Q$, and completeness holds for
every integer input and every odd $L\ge3$.

For every rational input, $\rho_{*,L}$ is a nonzero norm from $F_*$.
For each integer input, the source's prepared witnesses can be chosen so
that both contact weights $w_0,w_1$ are norms from $F_*$. Consequently
both conics
\begin{equation}\label{eq:bare}
                 u_j^2-d_*y_j^2=4w_j/17\qquad(j=0,1)
\end{equation}
have rational points.

No numerical sufficient value of $L$, degree bound for the resulting
defining polynomial, or realization with prescribed first coordinates
in~\eqref{eq:bare} is asserted.
\end{theorem}

\section{The four-contact norm identity}
Let
\[
 (\theta_0,\theta_1,\theta_2,\theta_3)=(0,1,2401,-343/32),
\]
\[
 (a_0',a_1',a_2',a_3')
 =(5951086353,5960872849,17757117551,5845158349).
\]
Direct integer identities give
\begin{equation}\label{eq:squares}
                         D_0^2+n\theta_i=(a_i')^2
                         \qquad(0\le i\le3).
\end{equation}

\begin{lemma}\label{lem:norm}
For any rational $Z$ and any nonsquare $d\in\Q^\times$, put
$\beta=(dZ^2-D_0^2)/n$. Then
\[
 r_4(\beta)=\frac{32}{n^4}
 \Nm_{\Q(\sqrt d)/\Q}\left(\prod_{i=0}^3(a_i'+Z\sqrt d)\right).
\]
\end{lemma}
\begin{proof}
Equation~\eqref{eq:squares} gives
$\beta-\theta_i=-\Nm(a_i'+Z\sqrt d)/n$.
Multiply the four identities, using
$r_4(\beta)=32\prod_i(\beta-\theta_i)$.
\end{proof}

The constants admit short certificates in $F_*$:
\begin{equation}\label{eq:constantnorms}
\begin{gathered}
 \nu_2=\frac{147+\sqrt{d_*}}{74},\quad \Nm(\nu_2)=2,\qquad
 \nu_{17}=\frac{143+\sqrt{d_*}}{24},\quad\Nm(\nu_{17})=17,\\
 \nu_\ell=\frac{1075327+3900\sqrt{d_*}}{3611},
 \qquad\Nm(\nu_\ell)=\ell.
\end{gathered}
\end{equation}
In particular $\Nm(4\nu_2)=32$.
The coefficient identities
\[
 A=d_*nH_0^2,\quad B=2d_*H_0K_0,\quad
 C=(d_*K_0^2-D_0^2)/n
\]
give
\begin{equation}\label{eq:z}
 Z(X)=nH_0U(X)+K_0,\qquad
 \beta_*(X)=\frac{d_*Z(X)^2-D_0^2}{n}.
\end{equation}
Thus $r_*(X)$ is an explicit norm. More precisely, if
\[
 \Gamma(X)=\frac{4\nu_2}{n^2}
             \prod_i(a_i'+Z(X)\sqrt{d_*}),\qquad
 \Gamma_C=\frac{4\nu_2}{n^2}
             \prod_i(a_i'+K_0\sqrt{d_*}),
\]
then $\Nm(\Gamma(X))=r_*(X)$ and $\Nm(\Gamma_C)=M_0$.
Consequently
\begin{equation}\label{eq:rhonorm}
 \rho_{*,L}(X)=
 \Nm\left(\frac{\eta\nu_\ell^L\Gamma_C\Gamma(X)}
                       {1+C_Lr_*(X)}\right).
\end{equation}

\section{The input polynomial and scalar preparation}
\begin{lemma}\label{lem:integrality}
The polynomial $\beta_*$ has integer coefficients and degree eight.
For every rational $X$,
\[
                 \beta_*(X)\in\Z\quad\Longleftrightarrow\quad X\in\Z.
\]
For every real $x$, $\beta_*(x)>C>2401$ and $r_*(x)\ge M_0>0$.
\end{lemma}
\begin{proof}
The divisibilities $576\mid A$ and $24\mid B$ give integer
coefficients in~\eqref{eq:input}. The leading coefficient is $A/576>0$.
The polynomial $U$ is integer-valued on $\Z$. Moreover,
\[
 24U(x)=(x^2-3x+1)^2+23\ge23,
\]
so the real assertions follow from positivity of $A,B$ and from the
four increasing positive factors of $r_4$ on $(2401,\infty)$.

Suppose $v_p(X)=-e<0$. Each factor $X-i$, $0\le i\le3$, has
valuation $-e$, and hence
\[
                        v_p(U(X))=-4e-v_p(24).
\]
The nonzero valuations of $A$ are $14$ at $2$, $6$ at $3$, $2$
at $17,\ell$, and $1$ at
\[
                 11,13,101,757,5939,7013,10657.
\]
The valuations of $B$ at $2,3,17,\ell$ are respectively $4,3,1,1$.
At the remaining listed primes they are nonnegative. If $p=2,e=1$,
the three terms $AU^2,BU,C$ have orders $0,-3,1$, so the middle
term uniquely dominates. If $p=2,e\ge2$, the leading term has
order $8-8e$, strictly less than $1-4e$ and $v_2(C)=1$.
At every other denominator prime the leading term uniquely dominates
and is negative: its order is $4-8e$ at $3$, $2-8e$ at $17,\ell$,
$1-8e$ at the remaining listed primes, and $-8e$ elsewhere.
Thus $\beta_*(X)$ is not integral. The converse is immediate.
\end{proof}

\begin{lemma}\label{lem:preparation}
For every real $X$,
\[
                0<\rho_{*,L}(X)<\eta^2/(A_0^2\ell)<1.
\]
Its valuation at $\ell$ is
\[
 v_\ell(\rho_{*,L}(X))=
 \begin{cases}
 L,&v_\ell(X)\ge0,\\
 L-8+32e,&v_\ell(X)=-e<0.
 \end{cases}
\]
For integer $X$ it has square reduced denominator, is a square in
$\Q_2,\Q_3,\Q_{17}$, and has exact valuations $(4,8,2)$ there.
\end{lemma}
\begin{proof}
The real bound follows from $r_*\ge M_0$ and
$C_L^2>A_0^2\ell^{L+1}$.
The exact residue certificates are
\[
 C\equiv10968,\quad M_0\equiv19666,\quad
 1+A_0M_0\equiv21661\pmod\ell.
\]
For $\ell$-integral $X$, $U(X)$ is integral and $\beta_*\equiv C$,
so both $r_*$ and the denominator are units. For
$v_\ell(X)=-e<0$, Lemma~\ref{lem:integrality} gives
$v_\ell(\beta_*)=2-8e$ and $v_\ell(r_*)=8-32e$.
The denominator has this latter valuation, proving the formula.

For integer $X$, $U(X)\in\Z$ and $\beta_*(X)-C$ has valuations
at least $4,3,1$ at $2,3,17$. The four contact-factor valuations at
the anchor $C$ are
\[
\begin{array}{c|rrrr}
 p&C&C-1&C-2401&32C+343\\\hline
 2&1&0&0&0\\
 3&0&2&1&1\\
 17&0&0&0&0
\end{array}
\]
Therefore $r_*(X)/M_0$ lies in $1+8\Z_2$, $1+3\Z_3$ and
$1+17\Z_{17}$, respectively. These groups consist of local squares.
The number $\ell$ is also a square at all three places, and
$A_0$ is divisible by $2,3,17$. Formula~\eqref{eq:scalar} gives the
local square assertions and the exact valuations $(4,8,2)$.

At an integer input the denominator $1+C_Lr_*$ is relatively prime
to $\eta\ell M_0r_*$. A prime dividing $\eta M_0$ divides $A_0$;
a prime dividing $r_*$ cannot divide $1+C_Lr_*$; and the
$\ell$ assertion follows from the residue certificate. Thus the
reduced denominator is a square.
\end{proof}

\section{Transfer through the general source}
Use $\beta_*,\rho_{*,L}$ in the source~\cite{UnitNorm}, with all
nine point pairs, sixteen comparisons, six scalar tests and local
orders unchanged. Retain its contact sets and poles
\[
 \Theta=\{0,-1/2,-1,1/2399,-32/407\},\quad
 \Theta_j=(4j+2+\Theta)\cup(4j+4+\Theta),\quad \gamma_j=-j-1,
\]
and put
\[
 a_0=1,\quad a_1=1+\rho_{*,L}T^2,\quad t_j=a_j\lambda,\quad
 \delta_j=t_j-\gamma_j,\quad
 R_j=\frac{\prod_{\theta\in\Theta_j}(t_j-\theta)}{\delta_j^{10}},
 \quad w_j=\rho_{*,L}R_j.
\]
The native products $\mathcal W_j$ and inputs
\begin{equation}\label{eq:native}
 b_j=\frac{4w_j}{\Tr(\mathcal W_j^{-1})^2}
\end{equation}
remain exactly those of that source.

\begin{proof}[Proof of the defining-family assertion in Theorem~\ref{thm:main}]
Let $S_*$ be $S$ together with the prime divisors of $M_0$.
It is fixed, and $\ell\notin S_*$ by Lemma~\ref{lem:preparation}.
Every finite odd contact of $\beta_*$ outside $S_*$ gives a positive
odd valuation of $\rho_{*,L}$: exactly one of the four factors in
$r_4$ vanishes, its coefficient factors and $M_0\eta\ell$ are units,
and the normalization denominator is a unit. There are no odd
denominator contacts outside $S$: the valuation formulas in
Lemma~\ref{lem:integrality} are even there. The fifth contact in the
inherited height estimate is infinity; no finite height contact has
been removed.

Omitting the finitely many prime divisors of $M_0$ changes only the
fixed height constant. The scalar-odd lemma of~\cite{General}, with
the unsquared-scale justification in~\cite{UnitNorm}, therefore gives
\[
               h(\beta_*)\le H_*c^{7/9},\qquad c\ge\ell^{2L+1}.
\]
Here $c$ is the nonzero adjacent root-index difference in the source.
The scalar has the uniform height bound
\[
               h(\rho_{*,L})\le12h(\beta_*)+2L\log\ell+O(1),
\]
since $r_4$ has degree four and all added coefficients are fixed.
All local comparison extraction and residual divisibilities are
unchanged. The aggregate trace and Lucas growth argument of the
unsquared source consequently has the same strict exponents, with
altered fixed constants, and forces $\beta_*(X)\in\Z$ for every
sufficiently large odd $L$. Lemma~\ref{lem:integrality} then gives
$X\in\Z$.

For completeness, an integer input gives a positive integer label
$\beta_*$ and a prepared scalar by Lemma~\ref{lem:preparation}.
The root choices, formal and real neighborhoods, free prime-pattern
slope and auxiliary-index refinements of the source apply in the
same order. The fixed prime divisors of $M_0$ may be included in the
permitted avoidance list before prime selection. Every inherited
point, denominator and whole-algebra guard is handled by the same
preparation and unit arguments. No new retained coordinate is used.
The finite-algebra encoding therefore has the same twenty-five total
variables. Its coefficient operations are finite and computable;
no previous numerical degree bound is transferred to this family.
\end{proof}

\section{Norms of the prepared weights}
\begin{lemma}\label{lem:prime}
Every rational prime $q\equiv1\pmod{10657}$, with $q\ne10657$,
is a norm from $F_*$.
\end{lemma}
\begin{proof}
The field is split at $2$ because $d_*\equiv1\pmod8$.
At $q$ it is split by quadratic reciprocity, and at $d_*$ the
number $q$ is a local square. At every other finite prime the
extension is unramified or split and $q$ is a unit, hence a norm.
The real extension is split. The conic Hasse principle
\cite[Theorem~26.1 and Corollary~26.5]{Poonen} gives a rational
point on $x^2-d_*y^2=qz^2$. Its coordinate $z$ is nonzero,
since $d_*$ is not a rational square.
\end{proof}

\begin{proof}[Proof of the norm-conic assertions in Theorem~\ref{thm:main}]
The scalar assertion is the identity~\eqref{eq:rhonorm}.
In the prime-pattern preparation include $d_*$ in the fixed odd
avoidance modulus. By~\cite[Section~3]{Dependencies}, every selected
Gaussian generator is then congruent to $i$ modulo $d_*$, so its
rational prime norm is congruent to one modulo $d_*$.
This is an allowed fixed avoidance condition and preserves the
required cross-order, denominator and real-approximation properties.

For a prepared slope $\lambda=u/v_0$, the exact source identity is
\[
 R_j=\left(\frac{a_j}{v_0\delta_j}\right)^{10}
       \prod_{\theta\in\Theta_j}
                     \bigl(u-(\theta/a_j)v_0\bigr).
\]
Every factor in the product is one of the selected norm primes by
Lemma~\ref{lem:prime}. The prefactor is a rational square and hence
a norm. Thus $R_j$ and $w_j$ are norms for both blocks. The numbers
$17$ and $4$ are norms by~\eqref{eq:constantnorms}, proving rational
solubility of~\eqref{eq:bare}.
\end{proof}

\begin{corollary}\label{cor:conditionalnorms}
Retain either conditional thirteen-variable source
of~\cite{Conditional,Direct}, including both direct-input height exponents
specified in~\cite{Direct,UnitNorm}, with its own arithmetic premises and
coefficient member. For every integer input in its point-stage
completeness statement, the two prepared physical weights can
simultaneously be chosen as norms from $F_*$. Both bare conics
\[
                       u_j^2-d_*y_j^2=4w_j/17,\qquad j=0,1,
\]
therefore have rational points. This requires no change to the source's
input, detector products, scales, physical weights or exceptional set.
\end{corollary}
\begin{proof}
Both literal conditional sources use
\[
 \Theta=\{0,-1/2,-1/4,-2/5,-1,-1/3\},\quad
 t_j=a_j\lambda,\quad \delta_j=t_j-1,\quad
 F_j=\prod_{\theta\in\Theta}(t_j-\theta).
\]
Their normalization is
\[
 b_j=\frac{4w_j}{\Tr(\mathcal W_j^{-1})^2},\qquad
 w_j=\frac{F_j}{4\delta_j^8}.
\]
For the original notation $R_j=F_j/\delta_j^8$, this is simply
$w_j=R_j/4$. Write the prepared common slope as $\lambda=u/v_0$.
The exact square-class identity is
\begin{equation}\label{eq:conditionalnormweight}
 w_j=
 \left(\frac{a_j^3}{2v_0^3\delta_j^4}\right)^2
 \prod_{\theta\in\Theta}
                  \bigl(u-(\theta/a_j)v_0\bigr).
\end{equation}
The rational-contact prime-pattern theorem makes the six contact
numerators selected positive rational primes, retaining any shared
incidences between the two blocks.

Include $d_*=10657$ in the permitted fixed avoidance modulus before
prime selection. It already belongs to the conditional exceptional
set $\{p:p\le76176\}$. The same Gaussian congruence used above
makes each selected prime congruent to one modulo $d_*$.
Lemma~\ref{lem:prime} makes every factor in
\eqref{eq:conditionalnormweight} a norm, and its prefactor is a square.
Thus both unchanged weights are norms. The allowed finite avoidance
retains the source's formal, real, denominator, reduction-order and
guard preparation. Dividing by the certified norm $17$ and multiplying
by the square $4$ proves the two conic assertions.
\end{proof}

For the finite-class reduction~\cite{FiniteClasses}, the needed
equations are more restrictive:
\[
 \bigl(m_j\Tr(\mathcal W_j^{-1})\bigr)^2-d_*y_j^2=4w_j/17
\]
with levels $m_j$ in a fixed finite set. Neither
Theorem~\ref{thm:main} nor Corollary~\ref{cor:conditionalnorms}
prescribes the first coordinates of its conic points.
Simultaneous rational realization of these actual trace coordinates,
with every inherited guard, remains the additional arithmetic premise
for the respective twenty-one- and nine-variable reductions. Bare
norm-conic solubility does not supply that premise.

\section{Exact certificates}
The standard-library program
\nolinkurl{computation/pole_parameters/norm_scalar_v47/check.py}
reconstructs the degree-eight polynomial, verifies the abstract norm
identity and all fixed norms, checks the uniform denominator-prime
inequalities and small-prime congruences, and authenticates the exact
source files. Its committed certificate includes the full input
coefficient list. The source-transfer proof and conic Hasse principle
are mathematical arguments and are not replaced by finite computation.

\begin{thebibliography}{9}
\bibitem{UnitNorm}
\emph{A single norm determinant for the 25- and 13-variable constructions},
project manuscript, 14 November 2026,
\nolinkurl{papers/finite_algebra/unit_norm_20261114_v28.tex}.
\bibitem{General}
\emph{A defining polynomial in twenty-five variables: combined scalar
and comparison conditions}, project source manuscript,
\nolinkurl{papers/few_variables25/few_variables25_merged_20261110_v19.tex}.
\bibitem{Conditional}
\emph{Thirteen Variables from a Weak $abc$ Inequality},
project source manuscript, 9 November 2026,
\nolinkurl{papers/conditional_few_variables/conditional13_20261109_v16.tex}.
\bibitem{Direct}
\emph{Thirteen Variables from Direct Elliptic Index Comparison},
project source manuscript, 10 November 2026,
\nolinkurl{papers/conditional_few_variables/direct_input13_20261110_v22.tex}.
\bibitem{Dependencies}
\emph{Arithmetic dependencies of the construction with 29 variables},
project manuscript, 8 November 2026,
\nolinkurl{papers/arithmetic_dependencies/arithmetic_dependencies_20261108_v8.tex}.
\bibitem{FiniteClasses}
\emph{Finite square-class conditions for native detector inputs},
project manuscript, 16 November 2026,
\nolinkurl{papers/pole_parameters/finite_complement_classes_20261116_v46.tex}.
\bibitem{Poonen}
B.~Poonen, \emph{Introduction to Arithmetic Geometry},
MIT 18.782 notes, 2009, Theorem~26.1 and Corollary~26.5,
\url{https://math.mit.edu/~poonen/782/782notes.pdf}.
\end{thebibliography}
\end{document}
